\section{Discussion}\label{sec:discussion}
\subsection{What the comparisons establish}
The central distinction is between following measured behavior and preserving a reference to approved behavior. Protected evolution retains deviations in this simulator, but the ordinary ARX already provides that benefit. The nonlinear correction and dual-window assembly do not meet the registered improvement target. Existing identification, residual correction, averaging and hysteresis remain close prior art; combining them does not by itself establish an original detector.

Reference accuracy, entry scoring and clearance logic have different roles. The no-dwell contrast shows that persistence can dominate benign occupancy despite unchanged scores and thresholds. AP hides alarm retention; schedule F1 hides much of physical recovery. Both must be accompanied by event and clearance measures.

\subsection{Relationship to published systems}
CASPER and its successor motivate external-sensor prediction; existing hybrid twins and identification methods already establish modeling and correction. E3 does not reproduce their full attack detectors. The Gaussian argument is standard error spending. Kaisar et al. already couple observer agreement with recovery transitions; their accepted manuscript now permits a method-level comparison, although its final publisher version and matching implementation remain unverified. Command-conditioned thresholds and motion-error intervals also predate our study.

MVT-Flow's offline likelihood, public fault labels and ordinary setpoints do not supply an online cyber-recovery benchmark or authenticated commands. Compatible full-method reproduction remains unresolved.

\subsection{Limitations and next evidence}
Three original parameterizations and three fresh perturbations share one simplified simulator. Noise seeds repeat stochastic realizations, not independent plants or acquisitions. Missing physics, trusted initialization and honest sensing limit external validity. The residual diagnostic receives known noise/bias/motion bounds; no operational procedure for estimating all of them is validated. Normal-only IMU development and metadata-only public-data admission do not provide a real-robot recovery test.

A worthwhile distinct mechanism remains open. Every completed condition is now exposed and cannot also confirm an outcome-informed replacement.